\section{Identification, Estimation, and Empirical Procedures}
\label{sec:methods}
\StageThreeSectionNote

This study is an observational measurement paper with disciplined theory rather
than a causal intervention study. The identified objects are the common-risk-set
margins \((q_h,\pi_h,\phi_h)\), the conditional timing summaries among observed
replies, and the thread-geometry outcomes defined in \Cref{sec:model}. Exposure
and staleness are not separately identified, and the policy indices in
\Cref{sec:model:or-diagnostic} are plug-in decision calculations on observed
margins rather than estimates of intervention effects.

\subsection{Archive discipline and current evidence boundary}
\label{sec:methods:design}

The canonical parent-linked archive is the public SimulaMet Moltbook snapshot
\citep{simulamet2026observatoryarchive}. The analysis begins with deterministic
deduplication, UTC normalization, and parent-link validation. Each non-root
comment becomes one candidate-parent unit, and each root post becomes one thread
for the geometry readouts. The present manuscript therefore reports only
SimulaMet-backed parent-linked evidence.

The study design also includes a separate larger-archive replication on a public
Moltbook archive (the approved MoltNet replication surface) under a harmonized
minimal schema. That replication is not yet implemented in the current evidence
base, so no pooled cross-archive estimator or cross-archive average is reported
here. Until minimal-schema equivalence is audited, archive disagreement is
treated as a replication object rather than pooled away.

\subsection{Primary risk-set estimators}
\label{sec:methods:def-est}
\label{sec:methods:two-part}

Let \(\mathcal{M}\) denote the realized cohort of non-root parent comments. For
parent \(m\in\mathcal{M}\) and horizon \(h\), define the realized within-horizon
completion indicator
\[
Y_h(m):=\mathbf{1}\{\delta_m=1,\ T_m\le h\}
\]
and the horizon risk-set indicator
\[
R_h(m):=\mathbf{1}\{C_m\ge h \text{ or } T_m\le h\}.
\]
The primary horizons are \(h\in\{5\text{ min},1\text{ h}\}\), with
\(30\) seconds used only in gap and censoring diagnostics.

The sample analogues of the control panel in \Cref{sec:model:estimands} are
\begin{align}
\widehat q_h
&=
\frac{\sum_{m\in\mathcal{M}} Y_h(m)}
     {\sum_{m\in\mathcal{M}} R_h(m)},
\\
\widehat \pi_h
&=
\frac{\sum_{m\in\mathcal{M}} \mathbf{1}\{\delta_m=1\}R_h(m)}
     {\sum_{m\in\mathcal{M}} R_h(m)},
\\
\widehat \phi_h
&=
\frac{\sum_{m\in\mathcal{M}} Y_h(m)}
     {\sum_{m\in\mathcal{M}} \mathbf{1}\{\delta_m=1\}R_h(m)}.
\end{align}
By construction, the sample identity
\[
\widehat q_h=\widehat \pi_h\,\widehat \phi_h
\]
holds exactly whenever the three quantities are computed on the same risk set.
Earlier draft tables denoted \(\widehat q_h\) by \(\widehat p_h\); the present
section uses \(\widehat q_h\) to reserve \(\widehat \pi_h\) and
\(\widehat \phi_h\) for its two interpretable margins.

The secondary descriptive incidence measure is
\[
\widehat p_{\mathrm{obs}}
=
\frac{1}{|\mathcal{M}|}\sum_{m\in\mathcal{M}} \mathbf{1}\{\delta_m=1\},
\]
the in-window ever-replied share. Because it is sensitive to heterogeneous
follow-up, \(\widehat p_{\mathrm{obs}}\) is treated as descriptive orientation
only and not as the paper's main decision object.

Among replied parents, conditional timing is summarized nonparametrically by the
empirical distribution of \(T_m\mid (\delta_m=1)\), especially
\((t_{50},t_{90},t_{95})\), plus short-window masses such as
\(\Pr(T_m\le 30\text{ s}\mid \delta_m=1)\) and
\(\Pr(T_m\le 5\text{ min}\mid \delta_m=1)\). Their unconditional counterparts
\(\Pr(\delta_m=1,\ T_m\le t)\) are reported on the same parent cohort.

\subsection{Gap handling, horizon standardization, and archive-conflict discipline}
\label{sec:methods:gap-discipline}

The canonical comment stream contains a 41.7-hour interruption, so gap handling
is part of the primary design rather than a late robustness appendix. The paper
therefore treats unobserved intervals as missing coverage rather than as true
silence and uses three built-in checks.

First, the core control-panel quantities are recomputed on contiguous windows,
including a short pre-gap window and the substantially longer post-gap window.
Second, parent comments whose potential follow-up windows overlap the gap are
excluded in prespecified overlap checks. Third, the headline completion
quantities are reported at explicit horizons using the risk-set estimators above,
so right-censoring near the observation boundary cannot masquerade as slow
conditional timing.

No replies are imputed across the gap. Likewise, if a later archive conflicts
with the canonical snapshot because of schema or coverage differences, that
conflict is handled through separate re-estimation rather than by pooled
cross-archive fusion. This is why gap handling, horizon standardization, and
archive-conflict discipline are treated as core identification design choices.

\subsection{Conversation geometry and downstream coordination readouts}
\label{sec:methods:geometry}

For each thread \(j\), depth is computed from \Cref{eq:depth}. The primary
geometry readouts are the empirical depth-tail probabilities
\(\Pr(D_j\ge k)\), maximum depth \(D_j\), and the branching-by-depth profile
\(\bar c_k\) from \Cref{eq:branching-by-depth}. We additionally summarize
reciprocity \(R_j\) and re-entry \(RE_j\) using
\Cref{eq:reciprocity-rate,eq:reentry-rate}. These outcomes are not treated as
independent mechanisms; instead, they provide the downstream coordination bridge
for testing whether lower \(q_h\) coincides with shallower continuation and
weaker repeated participation.

Where a compact depth-tail summary is useful, the paper reports an effective
log-tail slope \(\hat s_{\mathrm{depth}}\) fitted to
\(\log \Pr(D_j\ge k)\) for \(k\ge 2\). That slope is descriptive. It is not
claimed to be an identified reproduction parameter.

\subsection{Auxiliary stratified and parametric diagnostics}
\label{sec:methods:auxiliary}

Submolt heterogeneity is evaluated on the same risk-set estimands
\((\widehat q_h,\widehat \pi_h,\widehat \phi_h)\) using prespecified
deterministic topic strata. Topic heterogeneity is secondary; it remains useful
only insofar as it sharpens which persistence margin is more constrained.
Account-status contrasts are retained as descriptive heterogeneity only and do
not carry mechanism or causal interpretation.

For auxiliary descriptive incidence contrasts across coarse strata, we fit a
complementary log-log incidence regression:
\begin{equation}
\label{eq:methods-incidence-cloglog}
\log\!\left[-\log\!\left(1-\Pr(\delta_m=1\mid x_m)\right)\right]
=
x_m^{\top}\eta.
\end{equation}
Here \(x_m\) contains prespecified group indicators such as submolt class and
account-status group. This regression is calibration-oriented: it summarizes
direction and relative incidence differences, but it is not the paper's primary
estimand.

Parametric timing fits are also auxiliary only. For diagnostic comparison with
the nonparametric timing summaries, we fit an exponential-kernel hazard model
and a Weibull alternative. The exponential-kernel log-likelihood is
\begin{equation}
\label{eq:exponential-ll}
\ell(\alpha,\beta)
=
\sum_m
\left[
\delta_m(\log \alpha-\beta s_m)
-
\frac{\alpha}{\beta}\left(1-e^{-\beta s_m}\right)
\right],
\end{equation}
where \(s_m:=\min(T_m,C_m)\), and the Weibull survival benchmark is
\begin{equation}
\label{eq:weibull-survival}
S(s)=\exp\!\left[-\left(\frac{s}{\lambda}\right)^{\gamma}\right].
\end{equation}
These fits are used only as misspecification diagnostics and as secondary
timescale summaries; they do not determine the main conclusions.

\begin{remark}[Kernel half-life diagnostic]
\label{rem:estimand}
The quantity \(\ln 2/\beta\) from the exponential timing diagnostic is an
exponential-equivalent reply-hazard timescale. It is not a median thread
lifetime, and it is not a substitute for the primary control-panel margins.
\end{remark}

The repository analysis pipeline also reports a boundary-exclusion sensitivity
for the half-life diagnostic that removes parents created near the observation
end. That sensitivity is supplementary only and does not alter the primary
risk-set estimands.

\subsection{Diagnostic periodicity layer}
\label{sec:methods:periodicity}

Heartbeat-style synchronization is assessed only diagnostically and only on the
longest contiguous segment of the comment stream. The diagnostic stack includes:
event-time modulo-4-hour concentration, the Rayleigh statistic, Monte Carlo
\(p\)-values for exact phase uniformity, power spectral density (PSD) checks at
the four-hour target frequency, first-order autoregressive (AR(1)) null
calibration for binned-count spectra, and bin-width sensitivity.

This layer is intentionally non-load-bearing. At large \(N\), exact uniformity
can be rejected even when concentration is far too small to imply practically
meaningful global synchronization. Conversely, weak aggregate concentration does
not rule out individual routines under dephasing or jitter. Periodicity is
therefore treated as mechanism context only.

\subsection{Separate-archive replication requirement}
\label{sec:methods:replication}

The study design extends beyond the canonical first-week archive. The same
supported core metrics are to be re-estimated on a larger public Moltbook
archive under a harmonized minimal schema containing thread identifier, comment
identifier, parent identifier, author identifier, UTC timestamp, and community
label. Replication is performed separately by archive. No pooled cross-archive
estimator is used unless schema equivalence and coverage comparability have been
audited first.

If the replication archive preserves reliable parent-comment linkage, the same
geometry outcomes \((D_j,\bar c_k,R_j,RE_j)\) can be transported as well. If
parent linkage is incomplete or unreliable, the replication narrows to the
supported two-part persistence quantities \((q_h,\pi_h,\phi_h)\) and the
conditional timing summaries. In the current manuscript evidence base, however,
only the canonical SimulaMet archive supports the reported parent-linked
estimates.

\subsection{Comparator boundary and uncertainty}
\label{sec:methods:comparison}
\label{sec:methods:uncertainty}

No human-platform comparator enters the paper's primary estimands,
identification logic, or policy index. Any residual comparator analysis kept
elsewhere in the manuscript should be read as descriptive context only and not
as part of the flagship architecture.

Uncertainty for the nonparametric control-panel and geometry summaries is
reported with thread-cluster bootstrap intervals using fixed resampling rules.
Auxiliary incidence regressions use cluster-aware covariance that accounts for
dependence within threads and repeated authors where those identifiers are
available. One-parent-per-thread sensitivity is used as a further dependence
check. These procedures quantify uncertainty around observational readouts; they
do not convert the study into a causal design.
